%-*-latex-*-
%-*-latex-*-
\newcommand\COURSE{ciss240}
\newcommand\ASSESSMENT{q2001}
\newcommand\ASSESSMENTTYPE{Quiz}
\newcommand\POINTS{\textwhite{xxx/xxx}}

\input{myquizpreamble}
\input{yliow}
\input{\COURSE}
%-*-latex-*-

%-*-latex-*-

%-*-latex-*-
\renewcommand\TITLE{\ASSESSMENTTYPE \ \ASSESSMENT}

\renewcommand\EMAIL{}
\renewcommand\AUTHOR{eagraham1@cougars.ccis.edu}

\textwidth=6in
\begin{document}
\topmatterthree

This is a closed-book, no compiler, 5 minute quiz.

%------------------------------------------------------------------------------
\nextq
Write a function that prints the value in an integer array:
\begin{Verbatim}[frame=single,fontsize=\footnotesize]
int x[] = {1, 3, 5, 2, 4};
print(x, 3); // prints 1 3 5 (i.e., 3 values printed)
print(x, 4); // prints 1 3 5 2 (i.e., 4 values printed)
print(x, 5); // prints 1 3 5 2 4 (i.e., 5 values printed)
\end{Verbatim}
\ANSWER
\begin{answercode}
void print(int x[], int n)
{
    for (int i = 0; i < n; ++i)
    {
        std::cout << x[i] << " ";
    }
    std::cout << std::endl;
}
\end{answercode}

%------------------------------------------------------------------------------
\nextq
Complete the following program that returns the number of occurrences of the
value of \verb!target! in array \verb!x!.
The variable \verb!len! gives the number of values that the function should
test within array \verb!x!, i.e., the function search for occurrences of
\verb!target! in \verb!x[0]!, \verb!x[1]!, ..., \verb!x[len - 1]!.
For instance if \verb!y! is an array of doubles with values
$1.1$, $2.1$, $3.1$, $2.1$, $5.1$,
then calling \verb!count(y, 5, 2.1)! returns integer $2$ since \verb!2.1!
occurs twice in \verb!y!.
\\
\ANSWER
\begin{answercode}
int count(double x[], int len, double target)
{
    int occurrences = 0;
    for (int i = 0; i < len; ++i)
    {
        if (x[i] == target)
        {
            occurrences++;
        }
    }
    return occurrences;
}
\end{answercode}

%------------------------------------------------------------------------------
%-*-latex-*-
\newpage
\input{instructions}
\end{document}

